%% kapa: exam=2 round=21 year=2010 period=3 subject=compensation_law subject_ko=감정평가 및 보상법규 seq=1 number=1 points=40
%% source: Q-Net 감정평가사 기출문제 (cst00309) artlSeq=5024183 file=extracted/1008/5024183/2010년 제21회 감정평가사 2차시험 기출문제/3교시.hwp
%% conversion: kordoc 4.24.51 -> tools/gen2.py (auto); review=proofread-sonnet; flags=none
\begin{longquestion}{【문제 1】}
국토해양부장관은 전국을 철도로 90분 이내에 연결하기 위한 기본계획을 수립하였다. 이 계획에 기초하여 C공단 C이사장은 A지역과 B지역을 연결하는 철도 건설 사업에 대하여「공익사업을 위한 토지 등의 취득 및 보상에 관한 법률」(이하 “토지보상법”) 제20조에 따른 국토해양부장관의 사업인정을 받았다. P는 B-3공구 지역에 임야 3,000제곱미터를 소유하고 장뇌삼을 경작하고 있으며, 터널은 P 소유 임야의 한 가운데를 통과한다. C공단의 C이사장은 국토해양부장관이 제정한 K지침에 따라 P에 대하여 “구분지상권”에 해당하는 보상으로 900만원(제곱미터 당 3,000원 기준)의 보상금을 책정하고 협의를 요구하였다. P는 장뇌삼 경작임야에 터널이 건설되고 기차가 지나다닐 경우 농사가 불가능하다고 판단하여 C이사장의 협의를 거부하였다. (40점) \par
\medskip\noindent (1) P는 본인 소유 토지의 전체를 C이사장이 수용하여야 한다고 주장한다. 보상에 관한 C이사장의 결정과 P의 주장 내용의 정당성을 판단하시오. (20점) \par
\medskip\noindent (2) 토지보상법상 P가 주장할 수 있는 권리와 이를 관철시키기 위한 토지보상법상의 권리구제수단에 관하여 논술하시오. (20점) \par
\end{longquestion}
